\documentclass[a4paper, 12pt]{article}

\newcommand{\templates}{../../template}
\usepackage[a4paper, margin=2.5cm]{geometry}

\usepackage{enumitem}
\setlist[itemize]{noitemsep}
\setlist[enumerate]{noitemsep}

\let\oldpar\paragraph
\renewcommand{\paragraph}[1]{\oldpar{#1\\}\noindent}

% Avoid dots in the table of contents, it mess with the gulpease calculation
\makeatletter
\renewcommand{\@dotsep}{10000} 
\makeatother
\usepackage{graphicx}
\usepackage{hyperref}
\usepackage{makecell}
\usepackage{fancyhdr}

\newcommand{\settitolo}[1]{\newcommand{\titolo}{#1\\}}
\newcommand{\setprogetto}[1]{\newcommand{\progetto}{#1\\}}
\newcommand{\setcommittenti}[1]{\newcommand{\committenti}{#1\\}}
\newcommand{\setredattori}[1]{\newcommand{\redattori}{#1\\}}
\newcommand{\setrevisori}[1]{\newcommand{\revisori}{#1\\}}
\newcommand{\setresponsabili}[1]{\newcommand{\responsabili}{#1\\}}
\newcommand{\setversione}[1]{
	\ifdefined\versione\renewcommand{\versione}{#1\\}
	\else\newcommand{\versione}{#1\\}\fi
}
\newcommand{\setdestuso}[1]{\newcommand{\uso}{#1\\}}
\newcommand{\setdescrizione}[1]{\newcommand{\descrizione}{#1\\}}

\newcommand{\makefrontpage}{
	\begin{titlepage}
		\begin{center}

		\includegraphics[width=0.4\textwidth]{\templates/4ourSquared_logo}\\

		{\Large 4OURSQUARED}\\[6pt]
		\href{mailto://4oursquared.unipd@gmail.com}{4oursquared.unipd@gmail.com}\\
		
		\ifdefined\progetto
		\vspace{1cm}
		{\Large\progetto}
		{\large\committenti}
		\else\fi
		
		\vspace{1.5cm}
		{\LARGE\titolo}
		
		\vfill
		
		\begin{tabular}{r | l}
		\multicolumn{2}{c}{\textit{Informazioni}}\\
		\hline
		
		\ifdefined\redattori
			\textit{Redattori} &
			\makecell[l]{\redattori}\\
		\else\fi
		\ifdefined\revisori
			\textit{Revisori} &
			\makecell[l]{\revisori}\\
		\else\fi
		\ifdefined\responsabili
			\textit{Responsabili} &
			\makecell[l]{\responsabili}\\
		\else\fi
		
		\ifdefined\versione
			\textit{Versione} & \versione
		\else\fi
		
		\textit{Uso} & \uso
		
		\end{tabular}
		
		\vspace{2cm}
		
		\ifdefined\descrizione
		Descrizione
		\vspace{6pt}
		\hrule
		\descrizione
		\else\fi
		\end{center}
	\end{titlepage}
}
\usepackage{hyperref}
\usepackage{array}
\usepackage{tabularx}
\usepackage{adjustbox}

\newcounter{verscount}
\setcounter{verscount}{0}
\newcommand{\addversione}[5]{
	\ifdefined\setversione
		\setversione{#1}
	\else\fi
	\stepcounter{verscount}
	\expandafter\newcommand%
		\csname ver\theverscount \endcsname{#1&#2&#3&#4&#5}
}

\newcommand{\listversioni}{
	\ifnum\value{verscount}>1
		\csname ver\theverscount \endcsname
		\addtocounter{verscount}{-1}
		\\\hline
		\listversioni
	\else
		\csname ver\theverscount \endcsname\\\hline
	\fi
}

\newcommand{\makeversioni}{
	\begin{center}
		\begin{tabularx}{\textwidth}{|c|c|c|c|X|}
		\hline
		\textbf{Versione} & \textbf{Data} & \textbf{Redattore} & \textbf{Verificatore} & \textbf{Descrizione} \\
		\hline
		\listversioni
		\end{tabularx}
	\end{center}
	\clearpage
}

\settitolo{Norme di Progetto}
\setprogetto{Lumos Minima}
\setcommittenti{Imola Informatica}
\setredattori{Alberti Nicolas \\ Brotto Romina \\ Cavaliere Erica \\  Ceccato Francesco \\ Soldà Matteo}
\setdestuso{interno}
\setdescrizione{
Questo documento contiene le procedure, gli strumenti e i criteri di qualità che verranno usati nel progetto.
}

\addversione{0.0.0}{24/03/2023}{Soldà Matteo}{Brotto Romina}{Prima stesura.}
\addversione{0.0.1}{19/04/2023}{Soldà Matteo}{Ceccato Francesco}{Terminata Sezione Documentazione}
\addversione{0.0.2}{20/04/2023}{Ceccato Francesco}{Soldà Matteo}{Terminata Sezione Convenzioni di Codifica}
\addversione{0.0.3}{22/04/2023}{Alberti Nicolas}{Salami Lorenzo}{Stesura iniziale sezioni "1.1, 2.1, 2.2"}
\addversione{0.0.4}{22/04/2023}{Brotto Romina}{Cavaliere Erica}{Stesura iniziale sezioni "3.3, 3.4, 4, 4.2"}
\addversione{0.0.5}{24/04/2023}{Soldà Matteo}{Alberti Nicolas}{Separazione capitoli e risoluzione di alcuni typo}
\addversione{0.0.6}{24/04/2023}{Brotto Romina}{Soldà Matteo}{Risoluzione typo}
\addversione{0.0.7}{25/04/2023}{Alberti Nicolas}{Soldà Matteo}{Modifica elementi TODO e TBD; introduzione Glossario v.0.0.0}
\addversione{0.0.8}{26/04/2023}{Brotto Romina}{Alberti Nicolas}{Modifica elementi sezione 2.2.5}
\addversione{0.0.9}{03/05/2023}{Alberti Nicolas}{Soldà Matteo}{Aggiunta sezione 4.1.1; modifica sezione 4.1}
\addversione{0.0.10}{11/05/2023}{Erica Cavaliere}{Lorenzo Salami}{Aggiunta sezione 2.2.6}
\addversione{0.0.11}{24/07/2023}{Soldà Matteo}{Ceccato Francesco}{Inserite norme di codifica}
\addversione{0.1.0}{24/07/2023}{Ceccato Francesco}{Brotto Romina}{Verifica documento per RTB}
\addversione{1.0.0}{24/07/2023}{Soldà Matteo}{Ceccato Francesco}{Approvazione documento per RTB}

\begin{document}
\makeindexdetails
\makefrontpage \makeversioni
\tableofcontents
\newpage
\clearpage
\makecontentsdetails{Norme di Progetto} 

\section{Introduzione}
\subsection{Scopo del Documento}
Il documento si prefigge lo scopo di definire i metodi e le attività chiave
legate al \textit{Way of Working}\textsubscript{\textit{G}}. Tutti i membri del gruppo si impegnano a
seguire quanto riportato, al fine di ottenere consistenza e coerenza in tutti i
documenti e in tutte le attività del gruppo, per raggiungere in modo efficace ed
efficiente la realizzazione del prodotto finale.

\subsubsection{Scopo del prodotto}
Lo scopo del prodotto è quello di realizzare una \textit{Web Application}\textsubscript{\textit{G}} per
gestire e regolare l'intensità di uno o più impianti di illuminazione pubblica
in maniera automatica o manuale.

\subsubsection{Glossario}
Al fine di evitare ambiguità nei termini utilizzati durante la redazione dei
documenti, viene fornito anche il \textit{Glossario v.1.0.0,} dove vengono
definiti tutti i termini aventi un particolare significato. \\
Un termine presente nel \textit{Glossario} viene contrassegnato dal corsivo e
da una 'G' aggiunta in pedice.
\newpage

\section{Processi Primari}
\subsection{Fornitura}
\subsubsection{Scopo}
Lo scopo del \textit{processo}\textsubscript{\textit{G}} di fornitura è quello di descrivere e determinare ogni
compito ed attività svolta dal fornitore, al fine di comprendere e soddisfare le
richieste del proponente. 
La suddivisione e le tempistiche delle attività da svolgere saranno definite dal documento
\textit{Piano di Progetto}\textsubscript{\textit{G}}.

\subsubsection{Aspettative}
Le aspettative determinate dal processo di fornitura sono le seguenti:
\begin{itemize}
    \item Individuare e definire i bisogni che il prodotto deve soddisfare;
    \item Definire i requisiti e i vincoli dei processi;
    \item Stimare i costi per la realizzazione del prodotto;
    \item Ottenere dei feedback da proponente e committenti riguardanti il
    lavoro svolto;
    \item Chiarire eventuali dubbi sorti durante il progetto;
\end{itemize}

\subsubsection{Documenti}
\paragraph{\textit{Piano di Progetto}}
\\
Il \textit{Piano di Progetto} è il documento che viene redatto e mantenuto
durante tutta la durata del progetto. Al suo interno comprende:
\begin{itemize}
    \item \textbf{Analisi dei rischi:} vengono analizzati i rischi che possono
    essere individuati durante il corso del progetto e vengono descritte le
    strategie messe in atto al fine di contenerne la gravità;
    \item \textbf{Pianificazione del lavoro;}
    \item \textbf{Preventivo e consuntivo delle ore e dei costi.}
\end{itemize}
\paragraph{\textit{Piano di Qualifica}}
\\
Al fine di produrre materiale di qualità, viene redatto il documento
\textit{Piano di Qualifica}\textsubscript{\textit{G}}, che descrive le norme riguardanti la qualità dei prodotti e dei
processi del gruppo. Al suo interno comprende:
\begin{itemize}
    \item \textbf{Qualità di processo;}
    \item \textbf{Qualità di prodotto;}
    \item \textbf{Specifiche dei test;}
    \item \textbf{Resoconto delle attività di verifica.}
\end{itemize}

\subsubsection{Strumenti}
Nel processo di fornitura, il gruppo utilizza i seguenti strumenti per produrre
i documenti precedentemente elencati:
\begin{itemize}
    \item \textbf{Google Calendar:} sistema di calendari, usato per inserire le
    scadenze e gli eventi significativi riguardanti il gruppo;
    \item \textbf{Google Presentazioni:} è un
    software facente parte del gruppo di applicazioni Google, permette la
    realizzazione delle presentazioni del gruppo.
    \item \textbf{Google Fogli:} è un software facente parte del gruppo di
    applicazioni Google, permette la realizzazione di fogli di calcolo in cui
    vengono effettuati i conteggi delle ore. 
    \item \textbf{Online Gantt:} è un programma per realizzare i \textit{diagrammi di
    Gantt}\textsubscript{\textit{G}}.
\end{itemize}

\subsubsection{Repository Pubblica}
Durante lo sviluppo del progetto, il codice e la documentazione saranno depositati in una repository pubblica. La parte di interesse del proponente e del committente riguarda la cartella \textit{public}. \\
La cartella sopra citata contiene la documentazione necessaria alle revisioni di periodo con il committente. Tale cartella rispecchia la struttura della cartella \textit{docs}, situata nella parte privata. \\
I file che prevedono \textit{versionamento\textsubscript{G}}, riporteranno al loro interno la numerazione relativa alla versione attuale e un riassunto delle versioni precedenti. \\ 
Ogni modifica della parte pubblica deve essere approvata, tramite \textit{pull request}, dal responsabile.

\subsubsection{Contatti}
Tutti i contatti ufficiali, sia con il proponente che con il committente, avverranno tramite la mail del gruppo \href{mailto:4oursquared.unipd@gmail.com}{4oursquared.unipd@gmail.com}.\\
Le mail per la presentazione ad una revisione di periodo necessita di una conferma unanime, valutata la preparazione. Questa mail deve contenere la lettera di presentazione con i link ai vari documenti aggiunti/aggiornati nella parte pubblica della repository.

\subsection{Sviluppo}
\subsubsection{Scopo}
Il processo di sviluppo ha come scopo quello di descrivere le attività e i task di
analisi, progettazione, codifica, test, installazione e accettazione riguardanti
il prodotto software in sviluppo.

\subsubsection{Aspettative}
Le aspettative prevedono:
\begin{itemize}
    \item Determinazione dei requisiti del prodotto;
    \item Determinazione dei vincoli tecnologici e di design;
    \item Determinazione degli obiettivi di sviluppo;
    \item Realizzazione del prodotto finale, superando tutti i test individuati
    e soddisfando i requisiti e le aspettative del proponente.
\end{itemize}

\subsubsection{\textit{Analisi dei Requisiti}}

\paragraph{Scopo}
Lo scopo dell'\textit{Analisi dei Requisiti} è quello di:
\begin{itemize}
    \item Identificare requisiti obbligatori ed auspicabili richiesti dal
    proponente, tramite lo studio del capitolato;
    \item Supportare l'attività di pianificazione, fornendo informazioni utili
    al computo della mole di lavoro;
    \item Supportare l'attività di verifica, facilitandone il tracciamento dei requisiti.
\end{itemize}

\paragraph{Aspettative}
L'attività si pone come risultato la creazione di un documento contenente tutti
i requisiti richiesti dal proponente, includendo anche il punto di vista
dell'utente che utilizzerà i prodotti. Sarà presente una sezione legata al
tracciamento dai requisiti ai casi d'uso e viceversa.

\paragraph{Casi d'Uso}
Un Caso d'Uso è l'insieme degli scenari che prevedono uno stesso obiettivo per
un utente. Vengono descritte le interazioni con il sistema da parte di uno o più
attori, senza fornire alcun dettaglio implementativo. È formato da:
\begin{itemize}
    \item \textbf{Diagramma \textit{UML}} (opzionale): indica la relazione con
    altri Casi d'Uso;
    \item \textit{\textbf{Intestazione}}, nel formato:
    \begin{center}
        \begin{BVerbatim}
UC [Codice].[Sottocaso] - [Titolo]
        \end{BVerbatim}   
    \end{center}
    dove:
    \begin{itemize}
        \item \Verb^UC ^indica "\textit{Use Case}", ovvero "Caso d'Uso";
        \item \Verb^[Codice] ^è il codice identificativo del Caso d'Uso, 
        \item \Verb^[Titolo] ^è il titolo del Caso d'Uso. 
    \end{itemize}
    \item \textbf{\textit{Attore/i} primario/i};
    \item \textbf{\textit{Attore/i} secondario/i} (opzionale);
    \item \textbf{Descrizione};
    \item \textbf{Scenario principale};
    \item \textbf{Estensioni};
    \item \textbf{Inclusioni};
    \item \textbf{Precondizioni};
    \item \textbf{Postcondizioni};
    \item \textbf{Generalizzazioni}.
\end{itemize}

\paragraph{Requisiti}
I requisiti ricavati dall'analisi approfondita del capitolato, dalle discussioni
con i componenti del gruppo e dal confronto con il proponente, vengono suddivisi
in tre tipologie:
\begin{itemize}
    \item \textbf{Funzionali:} descrivono le funzionalità e il comportamento
    che il prodotto deve presentare;
    \item \textbf{Qualità:} descrivono i vincoli sulla qualità del prodotto e
    dei suoi componenti;
    \item \textbf{Vincolo:} descrivono i vincoli legati all'implementazione del
    prodotto, come ad esempio sulla tecnologia da utilizzare.
\end{itemize}
Ogni requisito ha una determinata importanza, di seguito descritta:
\begin{itemize}
    \item \textbf{Obbligatorio}
    \item \textbf{Desiderabile}
    \item \textbf{Opzionale}
\end{itemize}
Ogni requisito viene presentato nel formato univoco descritto di seguito: 
\begin{center}
    \begin{BVerbatim}
R[Tipologia][Caso d'Uso Relativo].[Sottocaso] - [Importanza]
    \end{BVerbatim}   
\end{center}
Dove:
\begin{itemize}
    \item \Verb^R^: Acronimo di "Requisito";
    \item \Verb^[Tipologia]^: indica la tipologia del requisito tra le seguenti:
    \begin{itemize}
        \item \textbf{F:} Funzionale;
        \item \textbf{Q:} Qualità;
        \item \textbf{V:} Vincolo;
    \end{itemize}
    \item \Verb^[Importanza]^: indica l'importanza del requisito tra le seguenti:
    \begin{itemize}
        \item \textbf{O:} Obbligatorio;
        \item \textbf{D:} Desiderabile;
        \item \textbf{P:} Opzionale.
    \end{itemize}
\end{itemize}

\subsubsection{Organizzazione dei File}
Tutta la documentazione sarà contenuta all'interno della cartella \texttt{docs/}. Questa sarà inoltre suddivisa in \textit{interna} ed \textit{esterna}, le quali conterranno rispettivamente la documentazione interna al gruppo e quella da condividere con il proponente e con il committente. \\
I verbali interni saranno conservati nella sottocartella \texttt{docs/interni/verbali/} mentre quelli esterni saranno conservati in \texttt{docs/esterni/verbali/}.\\
Per ogni documento esisterà una nuova sottocartella omonima.

\subsubsection{Convenzioni di Codifica}
\textbf{Convenzioni Linguistiche} \newline
\begin{itemize}
    \item Ogni elemento del codice deve essere scritto in lingua inglese, ad eccezione dei commenti, che devono essere scritti in italiano per evitare incomprensioni.
\end{itemize}
\paragraph{}\\
\textbf{Convenzioni di Nomenclatura File}
\begin{itemize}
    \item Un file che contiene l'\textit{interfaccia \textsubscript{G}} o l'\textit{implementazione\textsubscript{G}} di una sola ed esclusiva classe (ad eccezione di classi interne e sottoclassi) deve essere individuato nel seguente modo: \texttt{NomeClasse.estensione}.
\end{itemize}
\paragraph{} \\
\textbf{Convenzioni Stilistiche} \\
Per quanto riguarda la scrittura del codice, si è scelto di utilizzare le seguenti convenzioni, che sono state selezionate utilizzando lo \textit{standard de-facto} per il linguaggio \textit{Typescript}: \textit{camel case}, \textit{pascal case}, \textit{snake case} e \textit{screaming snake case}.\\ \\

\begin{center}
\begin{tabular}{|| c || c ||}
    \hline
        & TypeScript \\
    \hline \hline
    Funzioni e Metodi & camelCase() \\
    \hline
    Classi & PascalCase \\
    \hline
    Interfacce & PascalCase \\
    \hline
    Namespace e Package & camelCase \\
    \hline
    Costanti & SCREAMING\_SNAKE\_CASE \\
    \hline  
    Variabili Locali e Attributi & snake\_case \\
    \hline
    Variabili Globali & snake\_case \\
    \hline
\end{tabular}
\end{center}
\vspace*{0.5cm}
Per quanto riguarda i linguaggi di markup, i tag e gli attributi andranno scritti in lowercase se primitivi o in kebab-case se custom o complessi. \newline
Inoltre:
\begin{itemize}
    \item I membri di una classe dovrebbero essere dichiarati e/o definiti nel seguente ordine:
    \begin{itemize}
        \item Variabili di classe (statiche)
        \item Variabili di istanza (attributi) \\
        Nel seguente ordine (ove possibile):
        \begin{enumerate}
            \item public;
            \item protected;
            \item private;
        \end{enumerate}
        \item Costruttore;
        \item Distruttore (se richiesto);
        \item Metodi.
    \end{itemize}
\end{itemize}
\textbf{CSS}
\begin{itemize}
    \item Mantenere un file chiamato \texttt{style.css} per la versione desktop della webapp;
    \item Mantenere un file chiamato \texttt{global.css} per le regole che si applicano ad ogni aspetto del sito, in particolare ai moduli universali quali, a titolo esemplificativo e non esaustivo, header e footer;
    \item Mantenere un file chiamato \texttt{mini.css} con accorgimenti per la versione mobile;
    \item Specificare le unità di misura in \textit{em}, \textit{rem}, \textit{\%}, \textit{vw}, \textit{vh} (unità relative), a parte per i valori delle seguenti proprietà, i quali possono essere espressi in \textit{px}:
    \begin{itemize}
        \item \textit{text-shadow};
        \item \textit{box-shadow}.
    \end{itemize}
\end{itemize}

\paragraph{Indentazione e Blocchi di Codice}
\begin{itemize}
    \item Un \textit{blocco di codice\textsubscript{G}} delimitato da parentesi graffe dovrebbe essere strutturato apponendo le parentesi graffe ognuna in una nuova riga, non seguita da caratteri ulteriori;
    \item Si predilige l'uso della tabulazione per indentare la porzione di codice inclusa in un blocco.
\end{itemize}

\subsubsection{Workflow Github}

Ogni modifica alla repository dovrà seguire i seguenti passaggi:
\begin{itemize}
\item Modifica del codice su \textit{repository\textsubscript{G}} locale;
\item Push su un \textit{branch\textsubscript{G}} remoto;
\item Pull request tramite piattaforma web GitHub per merge sul branch \texttt{main};
\item Review della pull request;
\item Merge sul branch \texttt{main}.
\end{itemize}

Un branch remoto è un branch su cui operano uno o più membri del gruppo e ha come scopo l'elaborazione di un prodotto o di un processo di progetto. Il branch si chiamerà come il prodotto o il processo relativo.\newline
Le modifiche su un branch remoto avverranno tramite merge di un branch locale o con commit diretto. I membri del gruppo che lavorano sul branch concordano una modalità.\newline
È auspicabile che i branch locali non compaiano nella repository remota.\newline
Prima di effettuare un push alla repository remota o un merge è necessario fare uno squash dei commit superflui.\newline
Non appena chi lavora su un branch remoto ritenga che ci sia materiale pronto per una revisione, apre una pull request su GitHub.\newline
La review dovrà essere fatta dal verificatore. Il verificatore potrà accettare la pull request.\newline
Una pull request può essere accettata solo se tutti i prodotti coinvolti superano i requisiti minimi definiti nel \textit{Piano di Qualifica}.\newline
Quando il verificatore accetta un pull request, dovrà eliminare il branch relativo a quest'ultima.\newline
Non è consentito fare il caricamento diretto di file tramite interfaccia web o desktop.\newline
Non è consentito forzare i push sulla repository remota.\newline


\newpage

\section{Processi Secondari}
\subsection{Documentazione}
\subsubsection{Ciclo di Vita}
Ogni documento seguirà il seguente flusso:
\begin{itemize}
    \item \textbf{Pianificazione} \\ Il documento viene progettato ad alto livello secondo quanto richiesto;
    \item \textbf{Redazione} \\ Un membro del gruppo, nel ruolo di redattore, stila il documento;
    \item \textbf{Revisione} \\ Un membro del gruppo, nel ruolo di revisore, controlla l'assenza di errori grammaticali e il contenuto sia conforme alle norme di progetto;
    \item \textbf{Approvazione} \\ Un membro del gruppo, nel ruolo di responsabile, convalida il contenuto del documento corretto.
\end{itemize}

\subsubsection{Strumenti Utilizzati}
Per la stesura dei documenti verrà utilizzato il linguaggio \LaTeX. \\
Affinché i documenti risultino stilisticamente coerenti, dovranno includere al loro interno il file \texttt{docs/template/style.tex} per lo stile generale e il file \texttt{docs/template/front\_page.tex} per utilizzare il template della prima pagina.

\subsubsection{Struttura}
Ogni documento sarà così strutturato:
\begin{enumerate}
    \item Pagina di intestazione;
    \item Pagina con lista delle versioni, ove necessario;
    \item Pagina indice, ove necessario;
    \item Una o più pagine di contenuto.
\end{enumerate}

Le varie sezioni saranno così definite: \\

\textbf{Pagina di Intestazione} 
\begin{itemize}
    \item Logo e nome del gruppo;
    \item Nome del progetto e azienda cliente;
    \item Mail di contatto;
    \item Titolo del documento;
    \item Info generali:
    \begin{itemize}
        \item Redattori;
        \item Revisori;
        \item Responsabile;
        \item Versione, ove necessario;
        \item Destinazione d'uso (interno o esterno);
    \end{itemize}
    \item Sinossi.
\end{itemize}

\textbf{Elenco delle Versioni} \\
L'elenco delle versioni è una tabella così definita:
\begin{itemize}
    \item Versione;
    \item Data;
    \item Redattore;
    \item Revisore;
    \item Descrizione.
\end{itemize}
Le righe della tabella sono organizzate in ordine cronologico inverso. 
\newline \newline
\textbf{Indice} \\
L'indice è stato creato utilizzando il comando \textit{\textbackslash tableofcontents}
\newline \newline
\textbf{Norme di Progetto}
\begin{itemize}
    \item \textbf{Scopo:} Lo scopo del documento è quello di definire le procedure, gli strumenti e criteri di qualità al fine di stabilire un \textit{Way of Working}. Ogni membro del gruppo sarà tenuto a rispettare le indicazioni del documento presentate
    \item \textbf{Titolo:} Norme di Progetto
    \item \textbf{Nome del File:} \texttt{norme\_di\_progetto.tex}
\end{itemize} 
\paragraph{}\\
\textbf{Verbali}
\begin{itemize}
    \item \textbf{Scopo:} Lo scopo dei verbali è quello di rendicontare ciò che viene discusso durante una riunione, sia interna che esterna, delineando gli impegni che ne derivano
    \item \textbf{Titolo:} Ogni verbale sarà titolato con "$Verbale \ del \ <data> \ [con <esterni>]$"
    \item \textbf{Nome del File:} Tutti i file saranno chiamati \texttt{<yyyy>\_<mm>\_<dd>\_<tipo>.tex} dove $<yyyy>\_<mm>\_<dd>$ indica la data nella quale si è effettuata la riunione nel formato anno-mese-giorno; $<tipo>$ indica se la riunione è avvenuta tra i soli membri del gruppo "I" o con persone esterne "E".
    \item \textbf{Indice:} Non presente in quanto non utile
    \item \textbf{Struttura:}
    \begin{itemize}
        \item Data, ora e durata;
        \item Luogo;
        \item Partecipanti (interni ed esterni, raggruppati per appartenenza e ordinati alfabeticamente secondo il cognome);
        \item Ordine del giorno;
        \item Conclusioni derivanti dalla riunione;
        \item Impegni assunti.
    \end{itemize}
    \item \textbf{Stesura:} All'inizio di ogni riunione verrà nominato un membro del gruppo che redigerà il verbale e un membro che lo validerà. Chi redige il verbale avrà 24 ore per completare il lavoro, lo stesso tempo sarà concesso al validatore per approvarlo e segnalare le modifiche da applicare.
\end{itemize}
\paragraph{}\\
\textbf{Piano di Progetto}
\begin{itemize}
    \item \textbf{Scopo:} Lo scopo del documento è quello di definire gli obiettivi da raggiungere, stabilendo le tempistiche e il responsabile, tenendo traccia dei progressi e valutando i coste sostenuti rispetto a quelli preventivati.
    \item \textbf{Titolo:} Piano di Progetto
    \item \textbf{Nome del File:} \texttt{piano\_di\_progetto.tex}
\end{itemize}
\paragraph{}\\
\textbf{Piano di Qualifica}
\begin{itemize}
    \item \textbf{Scopo:} Lo scopo del documento è quello di definire le metriche e i requisiti minimi di qualità affinché un prodotto del progetto possa essere approvato.
    \item \textbf{Titolo:} Piano di Qualifica
    \item \textbf{Nome del File:} \texttt{piano\_di\_qualifica.tex}
\end{itemize}
\subsubsection{Convenzioni Documentali}
\begin{itemize}
    \item \textbf{Riferimento a file o cartelle:}
    \begin{itemize}
        \item Per fare riferimento al nome di un file o di una cartella si utilizza il \texttt{testo monospaziato};
        \item Il nome di una cartella deve sempre terminare con \texttt{/}.
    \end{itemize}
    \item \textbf{Stringhe e nomi:}
    \begin{itemize}
        \item Le stringhe vanno scritte tra doppi apici;
        \item Per indicare un parametro rappresentate una stringa si usa il testo tra parentesi angolari. Tale parametro non può contenere spazi;
        \item Per indicare una parte di nome o stringa opzionale si usa il testo racchiuso tra parentesi quadre.
    \end{itemize}
    \item \textbf{Riferimento tra documenti:}
    \begin{itemize}
        \item Riferimenti interni al documento: comando \texttt{$\backslash$ref};
        \item Riferimenti esterni al documento: nome della sezione del documento a cui si fa riferimento in corsivo.
    \end{itemize}
\end{itemize}
\subsection{Versionamento}
\subsubsection{Repository Github}
Per il versionamento del progetto si è deciso di utilizzare \textit{Git} sulla piattaforma \textit{GitHub}. La repository è raggiungibile tramite l'indirizzo \href{https://github.com/4ourSquared/LumosMinima}{https://github.com/4ourSquared/LumosMinima}.
\subsubsection{Convenzioni di Versionamento}
Per il versionamento si farà uso del Versionamento Semantico, secondo le \href{https://semver.org/lang/it/#specifica-di-versionamento-semantico-semver}{linee guida}.\newline
Questo tipo di versionamento si fonda sulle seguenti caratteristiche principali:
\begin{enumerate}
    \item Un numero di versione deve essere nella forma $X.Y.Z$ dove $X,Y,Z$ sono numeri interi non negativi e non devono contenere zeri iniziali;
    \item La versione \textit{major zero} (es: $0.X.Y$) serve solo per lo sviluppo iniziale e indica una versione non stabile del prodotto;
    \item La versione \textit{patch Z} (es: $x.y.Z$) deve essere incrementata solo se sono state introdotte correzioni retrocompatibili del testo (modifica da parte del redatore);
    \item La versione \textit{minor Y} (es: $x.Y.z$) deve essere incrementata se nella documentazione è stata introdotta una nuova modifica retrocompatibile (approvazione da parte del verificatore);
    \item La versione \textit{major X} (es: $X.y.z$) deve essere incrementata solo se nella documentazione è stata introdotta una modifica non retrocompatibile con le versioni precedenti (pubblicazione da parte del responsabile).
\end{enumerate}
\subsubsection{Comandi Utili}
\begin{itemize}
    \item \textbf{Sincronizzazione con la repository remota:} \texttt{git pull};
    \item \textbf{Creazione di un nuovo branch locale:} \texttt{git branch <nome\_branch>};
    \item \textbf{Passaggio ad un branch locale esistente:} \texttt{git checkout <nome\_branch>};
    \item \textbf{Aggiunta delle modifiche alla stage area:} \texttt{git add .};
    \item \textbf{Creazione del commit con le modifiche:} \texttt{git commit -m <descrizione>};
    \item \textbf{Push sul remote in un branch remoto non esistente:} \texttt{git push --set-upstream origin <nome\_branch>};
    \item \textbf{Push sul branch remoto esistente omonimo del branch locale:} \texttt{git push}.
\end{itemize}
Escludendo per il momento l'utilizzo di \textit{git flow}, questi risultano i comandi da utilizzare.

\subsection{Verifica}
\subsubsection{Regole Generali}
\begin{itemize}
    \item Il Verificatore su occupa di controllare le Pull Request aperte dagli altri membri del gruppo;
    \item Il Verificatore NON contribuisce ai prodotti che verifica;
    \item Il Verificatore NON modifica i prodotti che verifica;
    \item Il Verificatore, nel caso in cui trovi errori, è tenuto a commentare le PR usando la funzione apposita di GitHub.
\end{itemize}
\subsubsection{Metodi di Verifica}
Il verificatore può far uso dei test automatici presenti nel Piano di qualifica.\\
Per le sezioni di codice è consigliato un walkthrough del prodotto, similmente si dovrebbe chiedere a chi abbia scritto il codice eventuali delucidazioni su elementi poco chiari.

\subsection{Change Management}
Nel caso si dovessero cambiare sezione del codice bisogna controllare che eventuali modifiche
non vadano ad impattare altri elementi.\\
Per evitare ciò si farà uso del Diagramma delle classi presente nelle Specifiche architetturali.
\newpage
\section{Processi Organizzativi}
\subsection{Gestione dei Processi}
Si è scelto il modello Scrum con sprint di durata settimanale. Si è deciso di effettuare una riunione interna ogni lunedì dove discutere dei seguenti punti:
\begin{itemize}
    \item Resoconto dello Sprint precedente;
    \item Eventuali problemi riscontrati e soluzioni adottate;
    \item Resoconto di ore e costi per ogni componente;
    \item Pianificazione degli sprint successivi.
\end{itemize}
\paragraph{} I resoconti dovranno essere preparati dai singoli membri prima di ogni esposizione.\\ Ogni incontro dovrà produrre un aggiornamento delle milestones (e diagramma di Gantt). La pianificazione dei prossimi sprint consisterà nel dividere il lavoro in item e decidere il
numero di membri adatto per item. \\
\indent Gli item pianificati non potranno subire modifiche o cancellazioni durante il ciclo di sprint e dovranno essere portati a termine entro il ciclo successivo.
In caso emergessero importanti contestazioni o criticità il responsabile deciderà
come agire.\\
\indent Ogni ciclo di Scrum prevederà un nuovo responsabile, scelto per cognome in ordine alfabetico. Il responsabile dovrà assegnare le persone agli item considerando la
rotazione dei ruoli e assegnerà i membri alle attività previste tramite il
servizio di \textit{ticketing\textsubscript{G}}, come descritto nella sezione \nameref{sec:Ticketing}. \\
% \indent Le issues dovranno essere collegate alle milestones di appartenenza
% dal responsabile, che si occuperà di aggiornare preventivo e consultivo
% totale, compilando nel Piano di Progetto la sezione del ciclo di Scrum corrispondente. \\
% \indent Si è inoltre deciso che in caso di problemi singolarmente insormontabili andrà notificato
% il responsabile su Telegram che si occuperà di assegnare un membro per la risoluzione. \\
% \indent Il responsabile dovrà creare la issue su GitHub relativa al problema assegnandovici la persona
% scelta. Le soluzioni adottate per la risoluzione del problema andranno documentate
% dal responsabile stesso.
\subsubsection{Ticketing}
\label{sec:Ticketing}
Il ticketing è uno strumento che permette a tutti i membri del gruppo di
conoscere le attività da svolgere, quelle in corso e quelle completate. Il
Responsabile ne monitora l'avanzamento e può assegnare i compiti ai membri. Come
strumento di ticketing si utilizza il servizio di gestione delle issue integrato
in GitHub. Viene soddisfatto un requisito quando:
\begin{enumerate}
    \item viene creata una issue associata;
    \item viene collegata la issue ad un branch;
    \item i Verificatori devono controllare quanto è stato svolto: \begin{itemize}
        \item se la verifica ha esito positivo verranno effettuate le seguenti azioni: \begin{itemize}
            \item viene fatto il merge del branch nell'upstream;
            \item viene cancellato il branch creato appositamente per l'issue;
            \item l'issue viene chiusa;
            \item il requisito viene considerato soddisfatto. 
        \end{itemize}
        \item se la verifica ha esito negativo, verrà effettuata la seguente
        azione: \begin{itemize}
            \item vengono corretti gli errori segnalati dai Verificatori e si
            torna al punto (3).
        \end{itemize}
    \end{itemize}
\end{enumerate}
Le issue vengono create dal Responsabile o dai Verificatori e devono avere:
\begin{itemize}
    \item \textbf{title}: un titolo, dal contenuto breve e significativo;
    \item \textbf{description}: una descrizione, che deve essere breve, chiara e specifica, contenente il
    motivo dell'apertura della issue;
    \item \textbf{assignee}: il membro o i membri a cui viene associata l'issue.
    Possono variare durante il progetto;
    \item \textbf{\textit{milestone}}\textit{\textsubscript{G}}: la milestone ad
    essa associata;
    \item \textbf{labels}: le etichette, specificano lo stato della issue ed il
    contesto a cui è legata. Le etichette utilizzate sono le seguenti:
    \begin{itemize}
        \item \textbf{documentation:} indica che l'issue a cui si sta lavorando
        riguarda i documenti del gruppo.
        % AGGIUNGERE LABELS IN CASO DI UTILIZZO IN FUTURO
    \end{itemize}
\end{itemize}


\subsection{Miglioramento dei processi}
Durante l'attività di consuntivo di periodo dell’incontro settimanale, viene discusso il lavoro
fatto fino a quel momento, dando particolare peso ai problemi riscontrati. \\In base a ciò si valutano
le eventuali migliorie e correzioni applicabili.



\end{document}